% Figure 1: EML topology with two algorithms and null positions.
% 8-leaf binary tree. Algorithm A (blue) active from genesis.
% Algorithm B (amber) added at position 4; positions 0–3 are null (gray).
% Each node is a vertically-split rectangle: top = A, bottom = B.
%
% Designed for full-width (figure*) in USENIX two-column layout.

\begin{tikzpicture}[x=2.0cm, y=1.6cm]

% --- Shorthand ---
\def\hw{0.42}   % half-width of each node cell
\def\hh{0.18}   % half-height of each sub-cell (A or B)

% =====================================================================
% EDGES (drawn first so nodes sit on top)
% =====================================================================

% Level 0 → 1
\foreach \l/\r/\p in {0/1/0.5, 2/3/2.5, 4/5/4.5, 6/7/6.5} {
  \draw[black!30, line width=0.5pt] (\l, \hh) -- (\p, {1.0-\hh});
  \draw[black!30, line width=0.5pt] (\r, \hh) -- (\p, {1.0-\hh});
}
% Level 1 → 2
\foreach \l/\r/\p in {0.5/2.5/1.5, 4.5/6.5/5.5} {
  \draw[black!30, line width=0.5pt] (\l, {1.0+\hh}) -- (\p, {2.0-\hh});
  \draw[black!30, line width=0.5pt] (\r, {1.0+\hh}) -- (\p, {2.0-\hh});
}
% Level 2 → 3
\draw[black!30, line width=0.5pt] (1.5, {2.0+\hh}) -- (3.5, {3.0-\hh});
\draw[black!30, line width=0.5pt] (5.5, {2.0+\hh}) -- (3.5, {3.0-\hh});

% =====================================================================
% NODE DRAWING MACROS
% =====================================================================

% Active node: top=A (blue), bottom=B (amber)
\newcommand{\activenode}[2]{% #1=x, #2=y
  \fill[algA!15, draw=algA!65, line width=0.6pt]
    ({#1-\hw}, #2) rectangle ({#1+\hw}, {#2+\hh});
  \fill[algB!15, draw=algB!65, line width=0.6pt]
    ({#1-\hw}, {#2-\hh}) rectangle ({#1+\hw}, #2);
}

% Null node: top=A (blue), bottom=null (gray dashed + hatched)
\newcommand{\nullnode}[2]{% #1=x, #2=y
  \fill[algA!15, draw=algA!65, line width=0.6pt]
    ({#1-\hw}, #2) rectangle ({#1+\hw}, {#2+\hh});
  \fill[nullc!10, draw=nullc!50, line width=0.5pt, dashed]
    ({#1-\hw}, {#2-\hh}) rectangle ({#1+\hw}, #2);
  \fill[pattern=north east lines, pattern color=nullc!35]
    ({#1-\hw}, {#2-\hh}) rectangle ({#1+\hw}, #2);
}

% =====================================================================
% LEVEL 0: LEAVES
% =====================================================================

% Positions 0–3: A active, B null
\nullnode{0}{0}
\nullnode{1}{0}
\nullnode{2}{0}
\nullnode{3}{0}

% Positions 4–7: A active, B active
\activenode{4}{0}
\activenode{5}{0}
\activenode{6}{0}
\activenode{7}{0}

% =====================================================================
% LEVEL 1: INTERNAL
% =====================================================================

\nullnode{0.5}{1.0}
\nullnode{2.5}{1.0}
\activenode{4.5}{1.0}
\activenode{6.5}{1.0}

% =====================================================================
% LEVEL 2: INTERNAL
% =====================================================================

\nullnode{1.5}{2.0}
\activenode{5.5}{2.0}

% =====================================================================
% LEVEL 3: ROOT
% =====================================================================

% Root has a null left child and active right child for B.
% B still computes a valid root, so we show it as active.
\activenode{3.5}{3.0}

% =====================================================================
% LEAF INDEX LABELS
% =====================================================================

\foreach \i in {0,...,7} {
  \node[font=\small, text=black!60, anchor=north] at (\i, {-\hh-0.1}) {$d_{\i}$};
}

% =====================================================================
% EPOCH BOUNDARY (between positions 3 and 4)
% =====================================================================

\draw[densely dashed, red!55!black, line width=0.7pt]
  (3.5, -0.55) -- (3.5, {3.0+\hh+0.2});

\node[font=\footnotesize\itshape, text=red!55!black, anchor=south]
  at (3.5, {3.0+\hh+0.2}) {B added at position 4};

% =====================================================================
% ANNOTATION: representative hash labels
% =====================================================================

% Label on d₀ (null for B)
\draw[->, thin, black!50] (-0.9, -0.65) -- ({0-\hw+0.03}, {0-\hh+0.03});
\node[font=\footnotesize, anchor=east, text=black!70, align=right]
  at (-0.9, -0.65) {$H_a(d_0)$\\[-1pt]$N_0(b)$};

% Label on d₇ (active for both)
\draw[->, thin, black!50] (8.0, -0.65) -- ({7+\hw-0.03}, {0-\hh+0.03});
\node[font=\footnotesize, anchor=west, text=black!70, align=left]
  at (8.0, -0.65) {$H_a(d_7)$\\[-1pt]$H_b(d_7)$};

% =====================================================================
% LEGEND
% =====================================================================

\node[font=\footnotesize, anchor=center] at (3.5, -1.1) {%
  \raisebox{1pt}{\tikz \fill[algA!15, draw=algA!65, line width=0.5pt]
    (0,0) rectangle (0.3, 0.2);}
  ~Algo~A \qquad
  \raisebox{1pt}{\tikz \fill[algB!15, draw=algB!65, line width=0.5pt]
    (0,0) rectangle (0.3, 0.2);}
  ~Algo~B \qquad
  \raisebox{1pt}{\tikz{
    \fill[nullc!10, draw=nullc!50, line width=0.5pt, dashed]
      (0,0) rectangle (0.3, 0.2);
    \fill[pattern=north east lines, pattern color=nullc!35]
      (0,0) rectangle (0.3, 0.2);
  }}
  ~Null ($N_h$)%
};

\end{tikzpicture}
